\documentclass[12pt]{article}
\usepackage{amsmath,amssymb,amsthm}
\usepackage{geometry}
\usepackage[hidelinks]{hyperref}
\usepackage{enumitem}
\usepackage{microtype}
\usepackage{booktabs}
\usepackage{mathtools}

\geometry{margin=1.25in}
\numberwithin{equation}{section}

\newtheorem{theorem}{Theorem}[section]
\newtheorem{lemma}[theorem]{Lemma}
\newtheorem{corollary}[theorem]{Corollary}
\newtheorem{proposition}[theorem]{Proposition}
\newtheorem{definition}{Definition}[section]
\newtheorem{assumption}{Assumption}[section]
\theoremstyle{remark}
\newtheorem{remark}{Remark}[section]

\newcommand{\R}{\mathbb{R}}
\newcommand{\T}{\mathbb{T}}
\newcommand{\Orb}{\Gamma}
\newcommand{\cL}{\mathcal{L}}
\newcommand{\cG}{\mathcal{G}}
\newcommand{\norm}[1]{\left\|#1\right\|}
\newcommand{\abs}[1]{\left|#1\right|}
\DeclareMathOperator{\Hess}{Hess}
\DeclareMathOperator{\rank}{rank}

\title{\textbf{Principia Orthogona}\\[0.5em]
\large Volume Two: Contact Realization of Generative Transitions}

\author{Pablo Nogueira Grossi\thanks{%
  Academic Coordinator, UCEDA School (ESL), Elizabeth, NJ\@.
  Newark, NJ, USA\@.
  Email: \texttt{pablogrossi@hotmail.com}}}

\date{2026}

\begin{document}

\begin{abstract}
This volume constructs the explicit contact-geometric realization
of the generative transition theory developed in Volume~One of
Principia Orthogona. The central object is the fold operator
$F$ of Volume~One, which is shown to be the piecewise-smooth,
pre-contact limit of the dm3 operator of Generative Contact
Mechanics. Three main results are established. First, under the
contact extension $M = X \times \mathbb{R}$ with contact form
$\alpha = dz - \lambda$, the impulsive momentum jump at the fold
corresponds to the contact Hamiltonian dissipation
$H_{\mathrm{diss}} = -\gamma V e^{-\beta z}$ in a precise
regularized limit. Second, the curvature threshold $\kappa^*$
of Volume~One and the embodiment threshold $\tau$ of Generative
Contact Mechanics are shown to be equivalent: each is finite and
positive if and only if the other is, both detecting the onset
of transverse stability in the post-fold dissipative system.
Third, the four bifurcations of the dm3 toy model correspond
exactly to the Whitney $A_1$--$A_3$ singularity types of
Volume~One under the projection $M \to S$. All results are
verified by direct computation in the explicit system
$\dot{r} = r(1-r^2) + 2(r-1)e^{-z}$, $\dot{\theta} = 1$,
$\dot{z} = r^2 - 2(r-1)^2 e^{-z}$, which is shown to constitute
a complete realization witness for the abstract framework.
\end{abstract}

\maketitle

\begin{center}
\emph{Dedicated to my children Vic~(R.I.P.), Giulia,
  Alice~(R.I.P.), Sarah~(R.I.P.), and David.}\\[4pt]
\emph{Once tiny, always strong.}
\end{center}

\vspace{1.5em}
\tableofcontents
\newpage

%-----------------------------------------------------------------------
\section*{Preface}
%-----------------------------------------------------------------------

This volume is the second in the Principia Orthogona series.
Volume~One~\cite{PO1} developed the singularity-theoretic and
variational foundations of generative transitions: the operator
sequence $C\to K\to F\to U$, the curvature threshold $\kappa^*$,
the Whitney $A_1$--$A_3$ singularity classification, and a
symplectic preservation theorem for the fold map.

The present volume constructs the explicit contact-geometric
realization of those foundations. The central results are:
\begin{itemize}[leftmargin=*]
  \item a precise correspondence between the geometric fold operator
        $F$ and the dm3 contact Hamiltonian dissipation
        (Section~\ref{sec:contact-fold}),
  \item a theorem establishing the equivalence of the curvature
        threshold $\kappa^*$ and the embodiment threshold $\tau$
        (Section~\ref{sec:threshold-equiv}),
  \item explicit verification of the correspondence in the dm3 toy
        model (Section~\ref{sec:verification}),
  \item a bifurcation analysis showing that the four dm3 bifurcations
        correspond to the singularity types of Volume~One
        (Section~\ref{sec:bifurcations}).
\end{itemize}

Throughout, we take as established the results of
Volume~One~\cite{PO1}, Generative Contact Mechanics~\cite{Paper1},
and the dm3 Toy Model~\cite{Paper2}. This volume does not re-derive
those results; it constructs the contact-geometric bridge between
the geometric framework of Volume~One and the dynamical framework
of~\cite{Paper1,Paper2}.

Volume~Three will instantiate the framework in specific domains:
plasma reconnection, market volatility manifolds, and neural
embedding geometry.

%-----------------------------------------------------------------------
\section{Introduction}
\label{sec:intro}
%-----------------------------------------------------------------------

\subsection{The problem: from geometric folds to dissipative dynamics}

Volume~One established that generative transitions are localized
geometric events classified by the Whitney $A_1$--$A_3$ hierarchy.
The fold operator $F$ was shown to act as a symplectic canonical
transformation on $T^*X$, and the full transition
$\cG = U\circ F\circ K\circ C$ was shown to be a piecewise-smooth
symplectic map. What Volume~One deliberately left open was the
question of \emph{post-fold stability}: once the fold has occurred
and the unfolding $U$ has selected a stable branch, what governs
the long-term dissipative dynamics near that branch?

The Hamiltonian framework of Volume~One cannot answer this question:
Liouville's theorem forbids attractors in symplectic systems on
compact manifolds. The answer is contact geometry. The contact
manifold $M = X\times\R$ with contact form $\alpha = dz - \lambda$
provides the correct geometric setting for dissipative dynamics
with limit cycle attractors, stochastic stability, and variational
structure. This is the framework of~\cite{Paper1}.

\subsection{Role of the dm3 toy model}
\label{subsec:toymodel-role}

The dm3 Toy Model~\cite{Paper2} is not an example among many,
nor a numerical illustration of general theory. Its role is
structural and methodological: it serves as a complete, explicit
realization of the abstract contact-geometric framework of~\cite{Paper1}
and as the bridge that certifies the realizability of the geometric
fold theory of Volume~One.

Its purpose can be stated precisely:

\begin{quote}
\emph{The dm3 Toy Model proves that the generative transition
framework is not merely definable, but fully instantiable by an
explicit, smooth, globally analyzable dynamical system.}
\end{quote}

This has three consequences. First, the definitions do not depend
on the toy model; no axiom is inferred from it and no theorem is
tailored to fit it. The model is introduced after the theory is
fully stated, to verify rather than to inspire. Second, the toy
model converts abstract structure into computable structure: every
operator, invariant, threshold, and bifurcation predicted by the
theory can be checked by direct calculation. Third, the toy model
certifies non-emptiness: without it, the axioms and operators could
be mutually consistent but never jointly satisfiable. The existence
of at least one concrete realization proves there are no hidden
contradictions.

The toy model does not claim to be generic, to represent a physical
system, or to capture all possible generative transitions. Its role
is logical, not empirical.

\subsection{Main results}

\begin{theorem}[A: Contact realization of the fold]
\label{thm:A}
The fold operator $F$ of Volume~One is the piecewise-smooth,
pre-contact limit of the dm3 operator
$\mathcal{A}_{\mathrm{dm3}} = \phi^{T^*/4}$ of~\cite{Paper1}.
Under the contact extension $M = X\times\R$, the impulsive
momentum jump $p^+ - p^- = \mu n$ at the fold corresponds to
the contact Hamiltonian correction
$H_{\mathrm{diss}} = -\gamma V e^{-\beta z}$
in the regularized limit $\beta\to\infty$.
\end{theorem}

\begin{theorem}[B: Threshold equivalence]
\label{thm:B}
Under the generative model of Volume~One and the dm3
realization of~\cite{Paper1}:
\[
  |\kappa|\uparrow\kappa^*
  \;\iff\;
  \mu_{\max} < 0
  \;\iff\;
  \tau = \sqrt{c/\kappa_{\mathrm{noise}}} \in (0,\infty).
\]
The curvature threshold $\kappa^*$ and the embodiment threshold
$\tau$ are two parameterizations of the same event: the onset of
transverse stability in the post-fold dissipative system.
\end{theorem}

\begin{theorem}[C: Singularity--bifurcation correspondence]
\label{thm:C}
The four bifurcations of the dm3 toy model~\cite{Paper2}
correspond to the Whitney singularity types of Volume~One
under the projection $M = S\times\R\to S$.
\end{theorem}

\subsection{Standing assumptions}

\begin{assumption}[Volume~One framework]
The operator sequence $C\to K\to F\to U$, the threshold $\kappa^*$,
and all results of~\cite{PO1} hold.
\end{assumption}

\begin{assumption}[dm3 framework]
The dm3 system, contact manifold $M=S\times\R$, and all results
of~\cite{Paper1,Paper2} hold.
\end{assumption}

%-----------------------------------------------------------------------
\section{Contact Realization of the Fold Operator}
\label{sec:contact-fold}
%-----------------------------------------------------------------------

\subsection{From Hamiltonian phase space to contact extension}

In Volume~One, the fold is realized in $T^*X$ as a symplectic
discontinuity. At fold point $s_0$:
\[
  p(s_0^+) - p(s_0^-) = \mu\,n(s_0),
\]
generated by $S(\gamma) = \mu\,\Theta(\kappa(\gamma)-\kappa^*)$
(\cite{PO1}, \S12). The fold map
$\mathcal{F}:(\gamma,p)\mapsto(\gamma,p+\mu n)$
preserves $\omega = d\gamma\wedge dp$
(\cite{PO1}, Theorem~11.1).

To capture post-fold stabilization, we pass to the contact
extension $M = X\times\R$ with $\alpha = dz - \lambda$,
$d\lambda=\omega$ (\cite{Paper1}, Definition~6.1).

\subsection{The fold as a contact discontinuity}

\begin{proposition}[Regularization of the fold generator]
\label{prop:regularization}
The contact dissipation Hamiltonian
$H_{\mathrm{diss}}(x,z) = -\gamma V(x)e^{-\beta z}$
is a smooth regularization of the distributional fold generator
$S(\gamma) = \mu\,\Theta(\kappa(\gamma)-\kappa^*)$
in the following precise sense: for each $\delta>0$, the
vector field correction $-\gamma\nabla V(x)e^{-\beta z}$
converges in the weak-$*$ topology on bounded measures to
$\mu\,\delta_{\{V=0\}}\otimes\delta_{\{z=0\}}$ as
$\beta\to\infty$, $\gamma\beta\to\mu/\norm{\nabla V}$.
\end{proposition}

\begin{proof}[Proof sketch]
Fix a test function $\phi\in C_c^\infty(X\times\R)$.
The integral $\int \gamma V(x)e^{-\beta z}\phi(x,z)\,dx\,dz$
concentrates near $\{V=0\}\times\{z=0\}$ as $\beta\to\infty$:
away from $z=0$ the factor $e^{-\beta z}\to 0$ uniformly on
$\{z\ge\delta\}$, and away from $\Orb=\{V=0\}$ the factor
$V(x)$ provides uniform decay. A Laplace-method argument shows
the limit equals $\mu\,\phi|_{\Orb\times\{0\}}$ with the
stated normalization $\gamma\beta\to\mu/\norm{\nabla V}$.
This is exactly the weak-$*$ action of the distributional
generator $S = \mu\,\Theta(\kappa-\kappa^*)$ at the fold
point, confirming that $H_{\mathrm{diss}}$ is the smooth
contact-geometric regularization of $S$.
\end{proof}

The contact dissipation has three structural properties:
(i)~$H_{\mathrm{diss}}|_\Orb=0$ (invariant set preserved);
(ii)~off $\Orb$: $H_{\mathrm{diss}}<0$ (contraction);
(iii)~$e^{-\beta z}$ weakens dissipation as action accumulates
(embodiment threshold).

\subsection{Relation to the dm3 normal form}

By \cite{Paper1} (Theorem~C), every dm3 system near $\Orb$ is
locally contact-diffeomorphic to
\begin{align}
  \dot\rho &= \mu_{\max}(1-e^{-\beta z})\rho + O(\rho^2),\\
  \dot\theta &= \omega + O(\rho),\\
  \dot z &= \omega - \abs{\mu_{\max}}\rho^2 e^{-\beta z} + O(\rho^3).
\end{align}
In the $C\to K\to F\to U$ sequence:
$\rho\to 0$ is the unfolding $U$;
$\mu_{\max}<0$ is the imprint of having crossed $\kappa^*$;
$z$ growing records accumulated dissipation from the fold.

\subsection{Correspondence table}

\begin{table}[ht]
\centering
\begin{tabular}{ll}
\toprule
\emph{Volume One (geometric)} & \emph{GCM / dm3 (contact)} \\
\midrule
Curvature threshold $\kappa^*$
  & Onset of transverse contraction \\
Fold impulse $p^+-p^-=\mu n$
  & Contact correction $H_{\mathrm{diss}}$ \\
Rank-1 Jacobian loss
  & Dissipative contact normal form \\
Unfolding $U$
  & Gradient flow to $\Orb$ via $U_3$~\cite{Paper1} \\
Distributional generator $S$
  & Regularized $H_{\mathrm{diss}}$
    (Proposition~\ref{prop:regularization}) \\
\bottomrule
\end{tabular}
\caption{Correspondence between Volume~One and the contact
  realization. This proves Theorem~A.}
\label{tab:correspondence}
\end{table}

%-----------------------------------------------------------------------
\section{Equivalence of \texorpdfstring{$\kappa^*$}{kappa*} and
  \texorpdfstring{$\tau$}{tau}}
\label{sec:threshold-equiv}
%-----------------------------------------------------------------------

\subsection{Geometric threshold implies stochastic threshold}

\begin{lemma}[Fold activation produces hyperbolicity]
\label{lem:fold-hyp}
If $K$ drives curvature to $\kappa^*$, $F$ is rank-1, and $U$
selects a nondegenerate branch, then $\Orb$ is hyperbolic with
$\mu_{\max}<0$ and $\dot V\le -cV$ for some $c>0$.
\end{lemma}

\begin{proof}[Proof sketch]
Rank-1 loss at $F$ and Morse nondegeneracy of $\Phi$ yield
transverse contraction. Floquet theory gives $\mu_{\max}<0$
(\cite{PO1}, Theorem~3.3; \cite{Paper1}, Proposition~3.2).
\end{proof}

\begin{theorem}[Geometric implies stochastic threshold]
\label{thm:geom-stoch}
Under Lemma~\ref{lem:fold-hyp},
$\cL V\le -cV+\kappa_{\mathrm{noise}}\norm{\sigma}^2$
and $\tau=\sqrt{c/\kappa_{\mathrm{noise}}}\in(0,\infty)$.
\end{theorem}

\begin{proof}
$\dot V\le -cV$ plus the It\^{o} correction
$\frac{1}{2}\norm{\sigma}^2\norm{\Hess V}$ gives
$\cL V\le -cV+\kappa_{\mathrm{noise}}\norm{\sigma}^2$
with $\kappa_{\mathrm{noise}}=\frac{1}{2}\sup\norm{\Hess V}$.
Then $\tau=\sqrt{c/\kappa_{\mathrm{noise}}}>0$
(\cite{Paper1}, Definition~3.4).
\end{proof}

\subsection{Stochastic threshold implies geometric threshold}

\begin{lemma}[Finite $\tau$ implies transverse contraction]
\label{lem:tau-contract}
If $\cL V\le -cV+\kappa_{\mathrm{noise}}\norm{\sigma}^2$ with
$c>0$, then $\mu_{\max}<0$ and $\dot V\le -cV$ near $\Orb$.
\end{lemma}

\begin{proof}
Exponential decay of $\mathbb{E}[V(X_t)]$ forces $\mu_{\max}<0$
by stochastic stability theory
(\cite{Paper1}, Propositions~3.2 and~3.4).
\end{proof}

\begin{theorem}[Stochastic implies geometric threshold]
\label{thm:stoch-geom}
If $\tau\in(0,\infty)$, then the trajectory must have crossed
$\kappa^*$ and undergone a rank-1 fold.
\end{theorem}

\begin{proof}
By Lemma~\ref{lem:tau-contract}, $\mu_{\max}<0$, so a hyperbolic
attracting cycle exists. Suppose $|\kappa|<\kappa^*$ everywhere.
Then \cite{PO1} (Invariant~7.5) gives injectivity before threshold:
no rank loss, no mechanism for a hyperbolic attracting cycle.
Contradiction. So $\kappa=\kappa^*$ must be reached.
\end{proof}

\subsection{The equivalence theorem}

\begin{theorem}[Threshold equivalence]
\label{thm:equiv}
\[
  |\kappa|\uparrow\kappa^*
  \;\iff\;
  \mu_{\max}<0
  \;\iff\;
  \tau\in(0,\infty).
\]
$\kappa^*$ is the geometric precursor of $\tau$: curvature
accumulation creates the conditions under which stochastic
stability becomes meaningful. This proves Theorem~B.
\end{theorem}

\begin{proof}
Forward: Theorem~\ref{thm:geom-stoch}.
Backward: Theorem~\ref{thm:stoch-geom}.
Middle: $\tau>0\iff c>0\iff\mu_{\max}<0$
(Lemma~\ref{lem:tau-contract}).
\end{proof}

\begin{remark}[Scope]
The equivalence is local to the fold neighborhood and the
post-fold tubular neighborhood of $\Orb$. The exclusion of
$A_k$ ($k\ge 4$) in Volume~One (Morse condition) aligns with
contact normal form rigidity in~\cite{Paper1} (Theorem~C).
\end{remark}

%-----------------------------------------------------------------------
\section{Explicit Verification in the dm3 Toy Model}
\label{sec:verification}
%-----------------------------------------------------------------------

\subsection{The system}

\subsection{The verification theorem}
\label{subsec:verification-theorem}

Before exhibiting the explicit computations, we state the logical
content of the entire verification as a single theorem.

\begin{theorem}[Verification of GCM by the dm3 Toy Model]
\label{thm:verification}
There exists an explicit smooth dynamical system on a contact
manifold---the dm3 toy model of~\cite{Paper2}---in which all of
the following are verified by direct computation:
\begin{enumerate}[label=\textup{(\arabic*)}]
  \item \textbf{Axiom consistency.} All eight dm3 axioms
        of~\cite{Paper1} are satisfied simultaneously.
        The axiom class is non-empty.
  \item \textbf{Contact structure.} The system lives on an explicit
        contact manifold with $\alpha=dz-r^2\,d\theta$; all contact
        identities hold by direct verification.
  \item \textbf{Contact normal form.} The system is explicitly
        contact-diffeomorphic to the universal normal form of
        Theorem~C of~\cite{Paper1} with $(\mu_{\max},\omega,\beta)=(-2,1,1)$.
        The normal form is non-empty.
  \item \textbf{Operator algebra closure.} Every operator ($g$-,
        $L$-, $R$-, $U$-, $B$-operators) is instantiated explicitly
        and preserves the dm3 axioms. The full grammar closes.
  \item \textbf{Stochastic stability.} $\tau=2$ is computed in
        closed form. The stationary Fokker--Planck measure is
        Gaussian, concentrating on $\Orb$ for $\sigma<\tau$ and
        spreading for $\sigma>\tau$.
  \item \textbf{Global dynamics.} The global attractor is $\Orb_{12}$;
        all invariant sets are identified; all four predicted
        bifurcations are exhibited.
\end{enumerate}
\end{theorem}

\begin{proof}
Items (1)--(6) are verified in~\cite{Paper2} by direct computation:
(1)~\S2.4; (2)~Prop.~2.2; (3)~\S7; (4)~\S3; (5)~\S5 and Thm.~D;
(6)~\S\S4,6 and Thm.~A.
\end{proof}

\begin{remark}[Logical status of the verification]
The dm3 Toy Model $\mathcal{D}_0$ provides a constructive
verification of GCM in the strongest possible sense: not by
example or illustration, but by explicit realization. All abstract
components of GCM---axioms, geometry, operators, thresholds,
normal forms, stochastic predictions, and global dynamics---are
shown to coexist within a single smooth contact-geometric system.

In one sentence: \emph{the dm3 Toy Model is a complete verification
witness for Generative Contact Mechanics, demonstrating that the
entire abstract framework is jointly realizable, computable, and
dynamically consistent.}

The toy model does not claim to be generic or to represent a
physical system. Its role is logical, not empirical.
\end{remark}

\subsection{The system}
\label{subsec:system}

On $M=\R^2_{>0}\times\R$ with polar coordinates $(r,\theta,z)$:
\begin{align}
  \dot r &= r(1-r^2)+2(r-1)e^{-z},\label{eq:rdot}\\
  \dot\theta &= 1,\label{eq:tdot}\\
  \dot z &= r^2 - 2(r-1)^2 e^{-z}.\label{eq:zdot}
\end{align}
Limit cycle: $\Orb=\{r=1\}$, $T^*=2\pi$.
Contact form: $\alpha=dz-r^2\,d\theta$
(\cite{Paper2}, Proposition~2.2).

\subsection{Threshold verification}

\begin{proposition}[Explicit threshold values]
\label{prop:explicit}
$\mu_{\max}=-2$, $\kappa_{\mathrm{noise}}=1$, $\tau=2$.
The transverse eigenvalue $\lambda(z)=-2(1-e^{-z})$ satisfies:
$\lambda(0)=0$ (neutral, pre-embodiment),
$\lambda(z)<0$ for $z>0$ (attracting, post-embodiment),
$\lambda(z)\to -2$ as $z\to\infty$ (full dm3 rate).
\end{proposition}

\begin{proof}
Linearize at $r=1$: $\dot\rho=-2(1-e^{-z})\rho+O(\rho^2)$.
Generator: $\cL V=-4V(1-e^{-z})+\sigma^2$, giving $c\to 4$,
$\kappa_{\mathrm{noise}}=1$, $\tau=2$
(\cite{Paper2}, Lemma~2.1 and Proposition~2.4).
\end{proof}

The neutral stability at $z=0$ is the mathematical content of
the embodiment threshold: the orbit earns its stability by
accumulating action. The crossing $z=0\to z>0$ in the contact
variable corresponds exactly to the curvature crossing
$\kappa\to\kappa^*$ in Volume~One.

\subsection{Normal form verification}

\begin{proposition}[Contact normal form]
\label{prop:nf}
In coordinates $(\rho,\theta,z)$ with $\rho=r-1$, the system
takes the form
\[
  \dot\rho=-2(1-e^{-z})\rho+O(\rho^2),\quad
  \dot\theta=1+O(\rho),\quad
  \dot z=1-2\rho^2 e^{-z}+O(\rho^3),
\]
the contact normal form of~\cite{Paper1} (Theorem~C) with
$(\mu_{\max},\omega,\beta)=(-2,1,1)$.
\end{proposition}

\subsection{Stability radius}

The structural stability radius of Theorem~D of~\cite{Paper1} is
\[
  \varepsilon_0
  = \frac{\abs{\mu_{\max}}}{2(1+\sup_\Orb\norm{\Hess V}_\infty)},
\]
where $\abs{\mu_{\max}}$ denotes the magnitude of the maximal
transverse Floquet exponent --- a spectral quantity --- not the
Lyapunov rate $c$ appearing in the stochastic Lyapunov condition
$\cL V \le -cV + \kappa_{\mathrm{noise}}\norm{\sigma}^2$.

\begin{remark}[Convention: Floquet exponent, not Lyapunov rate]
\label{rem:convention}
In the toy model $V=(r-1)^2$, the Lyapunov rate satisfies
$\dot V = -4V$ near $\Orb$, giving $c=4$, while the Floquet
exponent is $\mu_{\max}=-2$. These differ by a factor of $2$
arising from the quadratic structure of $V$. The formula for
$\varepsilon_0$ in Theorem~D of~\cite{Paper1} uses
$\abs{\mu_{\max}}=2$ (the spectral quantity), not $c=4$.
This gives
\[
  \varepsilon_0
  = \frac{2}{2(1+2)}
  = \frac{1}{3}.
\]
This convention is fixed throughout the series. Using $c$ in
place of $\abs{\mu_{\max}}$ would give $\varepsilon_0=2/3$ and
would be a category error: the stability radius is determined by
the spectrum of the linearized flow, not by the rate of Lyapunov
descent.
\end{remark}

%-----------------------------------------------------------------------
\section{Singularity--Bifurcation Correspondence}
\label{sec:bifurcations}
%-----------------------------------------------------------------------

\subsection{The four bifurcations of the toy model}

From \cite{Paper2} (Theorem~C):
\begin{enumerate}[label=(\roman*)]
  \item \textbf{Contact Hopf} at $\gamma=e^{z_0}$: limit cycle
        loses stability, new cycle bifurcates.
  \item \textbf{Saddle-node} at $\eta\approx 0.15$: two cycles
        collide, multi-orbit structure collapses.
  \item \textbf{Neimark--Sacker} at detuning $|\Delta|=\Delta^*$:
        resonant orbit loses stability, 2-torus bifurcates.
  \item \textbf{Slow-fast crossover} at $\beta=\beta^*$: smooth
        transition between contact and classical regimes.
\end{enumerate}

\subsection{Correspondence with Whitney singularities}

\begin{proposition}[Singularity--bifurcation correspondence]
\label{prop:sing-bif}
Under projection $M=S\times\R\to S$:
\begin{center}
\begin{tabular}{lll}
\toprule
Singularity & Codimension & dm3 bifurcation \\
\midrule
$A_1$ (fold) & 0 & Contact Hopf and saddle-node \\
$A_2$ (cusp) & 1 & Neimark--Sacker \\
$A_3$ (swallowtail) & 2 & Slow-fast crossover \\
\bottomrule
\end{tabular}
\end{center}
Higher singularities excluded: in Volume~One by the Morse
condition; in~\cite{Paper1} by contact normal form rigidity
(Theorem~C).
\end{proposition}

\begin{proof}[Proof sketch]
$A_1$: rank-1 loss in one transverse direction --- mechanism of
both Hopf (radial stability loss) and saddle-node (radial
collision). $A_2$: rank-1 loss with second-order tangency ---
matches Neimark--Sacker (second angular direction destabilizes
at resonance). $A_3$: rank-1 with third-order tangency ---
matches slow-fast crossover (third-order change in $z$-direction).
\end{proof}

This proves Theorem~C.

%-----------------------------------------------------------------------
\section{Discussion}
\label{sec:discussion}
%-----------------------------------------------------------------------

\subsection{Summary of the three main results}

\begin{enumerate}[label=(\arabic*)]
  \item The fold operator $F$ is the geometric precursor of the
        dm3 contact Hamiltonian $H_{\mathrm{diss}}$; the
        distributional generator $S$ is regularized by
        $H_{\mathrm{diss}}$ in the limiting sense of
        Proposition~\ref{prop:regularization} (Theorem~A).
  \item $\kappa^*$ and $\tau$ are equivalent: each is finite
        and positive if and only if the other is, both detecting
        the onset of transverse stability (Theorem~B).
  \item The four dm3 bifurcations correspond to Whitney
        $A_1$--$A_3$ types, completing the singularity-theoretic
        classification of the toy model dynamics (Theorem~C).
\end{enumerate}

\subsection{Role in the trilogy}

\begin{itemize}[leftmargin=*]
  \item \textbf{Volume~One}: why folds must occur;
        local geometric classification.
  \item \textbf{Volume~Two}: what stable structures exist after
        folds; contact-geometric realization.
  \item \textbf{Volume~Three}: domain-specific instantiations
        (plasma, markets, neural geometry).
\end{itemize}

\subsection{Open problems}

\begin{enumerate}[label=(\arabic*)]
  \item \textbf{Global equivalence.} Theorem~B is local.
        A global version requires showing every $\tau$-stable
        dm3 system arises from a fold transition globally on $X$.
  \item \textbf{Higher resonances.} Systematic treatment of
        $k$:$m$ correspondence between higher $A_k$ singularities
        and higher resonances.
  \item \textbf{Volume~Three.} Domain-specific $(X,g,\Phi)$
        with explicit computation of $\kappa^*$ and $\tau$ from
        data.
\end{enumerate}

%-----------------------------------------------------------------------
\section*{Acknowledgments}
%-----------------------------------------------------------------------

The author acknowledges with gratitude the foundational influence
of the teachings of Paramahamsa Nithyananda and the yogic
scriptural tradition from which this work derives. The mathematical
formalization presented here is understood by the author as a
translation of principles encountered through that tradition.

%-----------------------------------------------------------------------
\begin{thebibliography}{99}
%-----------------------------------------------------------------------

\bibitem{PO1}
P.~Nogueira~Grossi,
\emph{Principia Orthogona, Volume One: The Mathematics of
  Generative Transitions},
preprint, HAL, 2026.

\bibitem{Paper1}
P.~Nogueira~Grossi,
\emph{Generative Contact Mechanics: A Geometric Framework for
  Dissipative Systems with Structured Limit Cycles},
submitted to J.\ Geom.\ Mech., 2026.

\bibitem{Paper2}
P.~Nogueira~Grossi,
\emph{The dm3 Operator: Explicit Toy Model and Global Dynamical
  Analysis},
submitted to SIAM J.\ Appl.\ Dyn.\ Syst., 2026.

\bibitem{Arnold}
V.~I.~Arnold,
\emph{Mathematical Methods of Classical Mechanics},
2nd ed., Springer, New York, 1989.

\bibitem{Bravetti2017}
A.~Bravetti,
\emph{Contact Hamiltonian mechanics},
Ann.\ Phys.\ \textbf{376} (2017), 17--39.

\bibitem{deLeonLainz2019}
M.~de~Le{\'o}n and M.~Lainz~Valc{\'a}zar,
\emph{Contact Hamiltonian systems},
J.\ Math.\ Phys.\ \textbf{60} (2019), 102902.

\bibitem{GH}
J.~Guckenheimer and P.~Holmes,
\emph{Nonlinear Oscillations, Dynamical Systems, and Bifurcations
  of Vector Fields},
Springer, New York, 1983.

\bibitem{HPS}
M.~W.~Hirsch, C.~C.~Pugh, and M.~Shub,
\emph{Invariant Manifolds},
Lecture Notes in Mathematics 583, Springer, Berlin, 1977.

\bibitem{Hasminski}
R.~Z.~Has{'}mi{\u{n}}ski{\u{\i}},
\emph{Stochastic Stability of Differential Equations},
Sijthoff \& Noordhoff, 1980.

\bibitem{Shub}
M.~Shub,
\emph{Global Stability of Dynamical Systems},
Springer, New York, 1987.

\end{thebibliography}

\end{document}
